% =============================================================
%  Planification des releases
% =============================================================

\begin{table}[H]
\centering
\caption{Planification des releases}
\label{tab:releases}
\small
\begin{tabular}{|c|p{3.5cm}|p{4.5cm}|c|p{4cm}|}
\hline
\textbf{Rel.} & \textbf{Titre} & \textbf{Objectif} & \textbf{Sprints} & \textbf{Livrables principaux} \\
\hline
R1 & Socle système et gestion hôtelière & Auth, sécurité, chambres, réservations, check-in QR & S1--S3 & Backend API, schéma DB, dashboards admin et réception, QR codes \\
\hline
R2 & Parcours client et services numériques & Espace chambre PWA, services hôteliers, messagerie bilingue + traduction IA & S4--S6 & Espace chambre PWA, services client, messagerie NLLB \\
\hline
R3 & Réclamations, interventions et intelligence intégrée & Réclamations + classification IA, Flutter worker, housekeeping & S7--S9 & Classification IA, dashboard manager, app Flutter, workflow housekeeping \\
\hline
\end{tabular}
\end{table}

\begin{quote}
\textit{Build, tests et validation} ne constituent pas une release séparée. Ils sont traités comme des activités continues (Definition of Done) et détaillés dans un chapitre technique séparé.
\end{quote}
